\begin{figure}[htbp]\centering
\begin{tikzpicture}[>=Latex]
\begin{scope}[shift={(0,0)}]
\draw[->,thick,cernblue] (0,0)--(2.5,0) node[right]{East};
\draw[->,thick,cernblue] (0,0)--(0,2.4) node[above]{North};
\draw[->,thick,cernblue] (0,0)--(-1.25,-.85) node[left]{Up};
\draw[->,teal,thick] (0,1.6) arc[start angle=90,end angle=0,radius=1.6];
\node[teal,font=\footnotesize] at (1.4,1.6) {azimuth};
\node[captionnode] at (.5,-1.6) {Navigation frame\\$E\times N=U$};
\end{scope}
\begin{scope}[shift={(7,0)}]
\draw[->,thick,cernblue] (0,0)--(2.5,0) node[right]{$+X$: west};
\draw[->,thick,cernblue] (0,0)--(0,2.4) node[above]{$+Y$: up};
\draw[->,thick,cernblue] (0,0)--(-1.5,-.9) node[left]{$+Z$: north};
\node[captionnode,text width=55mm] at (.4,-1.6) {Mount renderer coordinates\\East is $-X$, not $+X$};
\end{scope}
\end{tikzpicture}
\caption{The navigation and renderer bases are different but both preserve the intended chirality. The diagram is a projected axis sketch; up on the page is not a universal physical direction.}
\label{fig:frames}
\end{figure}
